\chapter{Tests}
Die Tests des Gesamtsystems teilen sich in folgende Gruppen auf:
\begin{itemize}
	\item \textbf{Datenbankzugriff:} Unter Anwendung des Testframeworks \enquote{JUnit} werden die Datenzugriffe -- also die SQL-Abfragen -- auf ihre Korrektheit geprüft. Darüberhinaus wird auch das Verhalten im Fehlerfall -- wie ungültige Anmeldedaten -- simuliert.
	\item \textbf{REST-API:} Mittels des in IntelliJ integrierten Http-Clients können definierte Http-Requests an die jeweiligen Endpoints gesendet werden; hiermit wird die fehlerfreie Funktion der REST-API sichergestellt. Zum Einen wird das Verhalten mit korrekten Eingabedaten (JSON-Bodies) zum Anderen das Erkennen von ungültigen bzw. unvollständigen Daten getestet.
	\item \textbf{App:} Zur Überprüfung des Verhaltens der App (Korrektes Navigieren, richtiges Verhalten der Buttons und fehlerfreies Kommunizieren mit dem Backend) werden manuelle Tests durchgeführt. Hierbei wird das alltägliche Benutzen der App nachgestellt und eventuell auftretende Fehler analysiert und gegebenenfalls beseitigt. Für die Studenten- und die Dozentenansicht wird jeweils ein Testaccount angelegt und die Anwendungsfälle dieser Rollen werden mehrmals pro Entwicklungsschritt zu Testzwecken ausgeführt.
\end{itemize}